\documentclass{ximera}
\title{Composition of Functions, Summary}

\newcommand{\pskip}{\vskip 0.1 in}

\begin{document}
\begin{abstract}
Composing functions.
\end{abstract}
\maketitle



\section{Solving Equations with $u$-Substitution}

Composing functions is a lot like composing actions. For example, consider the actions of putting on your socks and putting on your shoes. Two key points. First is that order matters. We get a completely different result if were to first put on our shoes and then our socks. And depending on how flexible your socks are, it might not even be possible to put your socks over your shoes.

The second point is that to undo the compostion of first putting on your socks and then your shoes, we undo the two actions in the \emph{reverse} order - we first take off our shoes and then our socks. And whether we're talking about socks and shoes or functions, it pays to remember this idea.

\begin{example} \label{E4433mzz}
Let
\[
  f(x)  = \frac{3x+5}{7x-8}
\]
and
\[
   g(x) = \frac{x-5}{1-x}.
\]

\begin{enumerate}
\item Evaluate $f(g(0))$.

\item Solve the equation 
\[
    f(g(x)) = 2.
\]
\end{enumerate}

\begin{explanation}
\begin{enumerate}
\item To evaluate $f(g(0))$, we work our way from the inside out and first evaluate $g(0)$. We get
\[
 g(0) = \frac{0-5}{1-0} = -5.
\]
We then input $g(0)$ into the function $f$. This gives
\begin{align*}
   f(g(0)) &= f(-5) \\
             &= \frac{3(-5) + 5}{7(-5) - 8} \\
              & = \frac{-10}{-43} .
\end{align*}
Our conclusion is that $f(g(0)) =10/43$. 

\item To solve the equation $f(g(x))=2$, we need to undo the action of the composition $f\circ g$. In evaluating $f(g(0))$, we first computed $f(0)$ and then $g(f(0))$. So to reverse the action of $f\circ g$ we first undo the action of $f$ and then the action of $g$. To undo $f$, we hide the inner function $g(x)$ by letting
\[
      u  = g(x).
\] 
Then the equation 
\[
  f(g(x)) =2
\]
becomes
\[
    f(u) = 2,
\]
or equivalently
\[
          \frac{3u+5}{7u-8} = 2.
\]
To solve this equation, multiply both sides by $7u-8$ to get
\[
 3u+5 = 2(7u-8),
\]
or equivalently
\[
3u + 5 = 14u - 16.
\]
Then 
\[
 21 = 11u
\]
and 
\[
 \frac{21}{11} = u.
\]

But remember that $u=g(x)$. So
\[
   g(x) = \frac{21}{11}
\]
and we now need to undo the action of $g$. For this, we substitute the given expression for $g(x)$ to get
\[
   \frac{x-5}{1-x} = \frac{21}{11}.
\]
Then multiply both sides of this equation by $11(1-x)$, giving
\[
  11(x-5) = 21(1-x).
\]
So
\[
 11x - 55 = 21 - 21x,
\]
\[
   32x = 76
\]
and 
\[
  x = \frac{19}{8} .
\]
Our conclusion is the solution set to the equation $f(g(x))= 2$ is
\[
    \{  x \, | \, x = \frac{19}{8} \}.
\]
\begin{question} \label{Q2321}
Check that $x=19/8$ is the solution to the equation $f(g(x))=2$ by evaluating $f(g(19/8))$.
\end{question}
\end{enumerate}

\end{explanation}
\end{example}


\section{Finding the Domain of a Composition}
\begin{example} \label{Edf300}
Let
\[
    f(x) = \frac{24}{x^2-x-6}
\]
and 
\[
    g(x) = \frac{3x-5}{x-4} .
\]
\begin{enumerate}
\item Evaluate $f(g(4))$.

\item Find the domain of the composition $f\circ g$.
\end{enumerate}

\begin{explanation}
\begin{enumerate}
\item Since 
\[
 g(4) = \frac{3(4) - 5}{4 -4} = \frac{7}{0},
\]
$f(g(4))$ is undefined.

That's an example of one thing that can go wrong in evaluating the composition $f(g(x))$; the input $x$ might not be in the domain of $g$. The other thing that can go wrong is that $g(x)$ might not be in the domain of $f$.

So the point is this. The number $x$ is in the domain of the composition $f\circ g$ if and only if
\begin{enumerate}
\item $x$ is in the domain of $g$ and

\item $g(x)$ is in the domain of $f$.
\end{enumerate}

So the first step in finding the domain of $f\circ g$ is to find the domain of $g$. For our particular example where
\[
     g(x) = \frac{3x-5}{x-4},  
\]
the only thing that can go wrong is that the denominator is zero. So $x$ is in the domain of $g$ if and only if
\[
   x - 4 \neq 0,
\]
or if and only if
\[
     x \neq 4.
\]

The next step is to make sure $u=g(x)$ is in the domain of $f$. To do this we first find the domain of
\[
      f(u) = \frac{24}{u^2-u-6}.
\]
Again, the only thing that can go wrong is if the denominator is zero. So $u$ is in the domain of $f$ if and only if
\[
    u^2 - u - 6 \neq 0,
\]
or equivalently if and only if
\[
  (u-3)(u+2) \neq 0.
\]
Now for the product of two numbers to be non-zero it is necessary and sufficient that each number be not equal to zero. So 
\[
  (u-3)(u+2) \neq 0
\]
if and only if
\[
       u - 3 \neq 0  \hskip 0.2 in \text{and} \hskip 0.2 in u+2 \neq 0.
\]
So
\[
   u \neq 3  \hskip 0.2 in \text{and} \hskip 0.2 in u \neq -2.
\]
 
So $u$ is in the domain of the function $f$ if and only if $u\neq -2$ and $u\neq 3$.

But remember for $x$ to be in the domain of $f(g(x))$, we need to make sure that $u=g(x)$ is in the domain of $f$. So $u=g(x)$ is in the domain of $f$ if and only if
\[
         g(x) \neq 3 \hskip 0.2 in \text{and} \hskip 0.2in  g(x) \neq -2,
\]
or if and only if
\[
   \frac{3x-5}{x-4} \neq 3 \hskip 0.2 in \text{and} \hskip 0.2in  \frac{3x-5}{x-4} \neq -2.
\]
Solving these inequalities gives
\[
 0x = 7  \hskip 0.2 in \text{and} \hskip 0.2in x\neq \frac{13}{5}.
\]

The equation $0x = 7$ has no solution and does not give an additional restriction on the domain of the composition $f(g(x))$. 

So remembering that $x$ must also be in the domain of $g$ (ie. $x\neq 4$), the domain of the composition $f(g(x))$ is the set
\[
 \{ x \, | \, x\neq 4  \text{ and }  x\neq 13/5 \}.
\]

\end{enumerate}
\end{explanation}

\end{example}


\begin{example}  \label{Erer54r56t364}
Use the graphs of the functions $y=f(x)$ and $y=g(x)$ below to find approximate answers to the following questions. Drag points $A$ and $B$ as needed to help with your approximations. Use set notation to write solution sets when solving an equation.

\begin{onlineOnly}
    \begin{center}
\desmos{9uxkofytwi}{900}{600}
\end{center}
\end{onlineOnly}

\href{https://www.desmos.com/calculator/9uxkofytwi}{141: Composition 2}

\end{example}


\end{document}